% Phần 4: Yêu cầu 3 (so sánh phương pháp, chi phí) và hạn chế.
% Mọi con số lấy qua \so{...} từ so_lieu.tex.

\section{Yêu cầu 3: so sánh phương pháp và chi phí}
\label{sec:yc3}

\subsection{Chọn mô hình: bốn ứng viên trên tập dev}
\label{sec:yc3-chon}

Nhóm huấn luyện song song bốn ứng viên trên cùng tập dữ liệu, cùng phân hoạch phân tầng theo nhóm trùng lặp
(\so{dev.chia.fit} mẫu \textbf{fit}, \so{dev.chia.dev} mẫu \textbf{dev}, \so{dev.chia.cal} mẫu \textbf{cal}), cùng
seed 42 và cùng tiêu chí chọn checkpoint. Ba ứng viên sử dụng chung kiến trúc bộ mã hóa \textbf{SafeVi-DualORC-CRF}
(phân loại 5 nhãn kết hợp gắn thẻ BIO + CRF mức từ), chỉ khác bộ mã hóa nền: mmBERT-base (đa ngôn ngữ), BamiBERT và
ViSoBERT (tiếng Việt). Ứng viên thứ tư (\textbf{GuardMerge}) chuyển đổi từ nền mô hình ngôn ngữ lớn hai chiều khoảng
0,6 tỉ tham số và tinh chỉnh bằng LoRA. Toàn bộ quy trình chọn mô hình được thực hiện độc lập trên tập \textbf{dev}
trước khi mở khóa tập kiểm thử (Hình~\ref{fig:p4-ung-vien}, Bảng~\ref{tab:p4-ung-vien}).

\begin{figure}[H]
\centering
% Sinh tự động bằng `python -m src.v2.eval.xuat_so_lieu_bao_cao`. Không sửa tay.
\begin{tikzpicture}
\begin{axis}[width=0.85\linewidth, height=6cm, xmin=0, xmax=100, ymin=89, ymax=93.6,
  xlabel={Độ trễ p50 mỗi đoạn trên GPU (ms, batch 1)}, ylabel={Macro-F1 trên dev},
  axis x line*=bottom, axis y line*=left, grid=major, grid style={draw=black!12},
  tick label style={font=\small}, label style={font=\small}, clip=false,
  yticklabel style={/pgf/number format/.cd, use comma, fixed}]
\addplot[only marks, mark=*, mark size=4pt, mauvhai] coordinates {(37.96,92.29)};
\node[anchor=west,xshift=4pt,font=\footnotesize] at (axis cs:37.96,92.29) {mmBERT-base (chọn): 92,3; 1.177 MiB};
\addplot[only marks, mark=*, mark size=4pt, navy!70] coordinates {(18.65,92.24)};
\node[anchor=north,yshift=-4pt,font=\footnotesize] at (axis cs:18.65,92.24) {BamiBERT: 92,2; 397 MiB};
\addplot[only marks, mark=*, mark size=4pt, navy!70] coordinates {(20.82,92.89)};
\node[anchor=south,yshift=4pt,font=\footnotesize] at (axis cs:20.82,92.89) {ViSoBERT: 92,9; 376 MiB};
\addplot[only marks, mark=*, mark size=4pt, navy!70] coordinates {(82.31,90.03)};
\node[anchor=south,yshift=4pt,font=\footnotesize] at (axis cs:82.31,90.03) {GuardMerge 0,6B: 90,0; 1.177 MiB};
\end{axis}
\end{tikzpicture}

\caption{Bốn ứng viên trên tập dev: Macro-F1 theo độ trễ GPU (batch~1), kèm dung lượng trọng số FP32. Ba bộ mã hóa
SafeVi-DualORC-CRF vượt trội GuardMerge cả về chất lượng lẫn tốc độ.}
\label{fig:p4-ung-vien}
\end{figure}

\begin{table}[H]
\centering
\small
\caption{So sánh bốn ứng viên trên tập dev (F1 trên thang 0--100; FPR và Tô oan tính theo \%). ``Tô oan'': tỉ lệ đoạn
\nhan{AN_TOAN} bị trích xuất bằng chứng. ``Chênh vùng'': hiệu Macro-F1 giữa vùng cao nhất và thấp nhất. Độ trễ đo trên
GPU \so{v2.gpu}, batch~1.}
\label{tab:p4-ung-vien}
\setlength{\tabcolsep}{4pt}
\begin{tabular}{lcccc}
\toprule
\textbf{Chỉ số} & \textbf{mmBERT-base} & \textbf{BamiBERT} & \textbf{ViSoBERT} & \textbf{GuardMerge} \\
\midrule
Macro-F1 (5 nhãn) & \so{dev.mmbert.macro} & \so{dev.bamibert.macro} & \textbf{\so{dev.visobert.macro}} & \so{dev.guard.macro} \\
Binary-F1 (nguy hại/an toàn) & \so{dev.mmbert.bin} & \so{dev.bamibert.bin} & \textbf{\so{dev.visobert.bin}} & \so{dev.guard.bin} \\
FPR đoạn \nhan{AN_TOAN} & \so{dev.mmbert.fpr} & \so{dev.bamibert.fpr} & \textbf{\so{dev.visobert.fpr}} & \so{dev.guard.fpr} \\
Bằng chứng: BC F1 theo ý & \so{dev.mmbert.bcf} & \textbf{\so{dev.bamibert.bcf}} & \so{dev.visobert.bcf} & \so{dev.guard.bcf} \\
Tô oan đoạn \nhan{AN_TOAN} & \so{dev.mmbert.toan} & \so{dev.bamibert.toan} & \textbf{\so{dev.visobert.toan}} & \so{dev.guard.toan} \\
Chênh vùng (Macro-F1) & \textbf{\so{dev.mmbert.gap}} & \so{dev.bamibert.gap} & \so{dev.visobert.gap} & \so{dev.guard.gap} \\
Macro-F1 tiếng lóng ($n = \so{dev.tieng.n}$) & \textbf{\so{dev.mmbert.tieng}} & \so{dev.bamibert.tieng} & \so{dev.visobert.tieng} & \so{dev.guard.tieng} \\
\midrule
Trọng số FP32 (MiB) & \so{dev.mmbert.trongso} & \so{dev.bamibert.trongso} & \textbf{\so{dev.visobert.trongso}} & \so{dev.guard.trongso} \\
Độ trễ p50 (ms) & \so{dev.mmbert.p50} & \textbf{\so{dev.bamibert.p50}} & \so{dev.visobert.p50} & \so{dev.guard.p50} \\
Thời gian huấn luyện (phút) & \so{dev.mmbert.trainphut} & \textbf{\so{dev.bamibert.trainphut}} & \textbf{\so{dev.visobert.trainphut}} & \so{dev.guard.trainphut} \\
\bottomrule
\end{tabular}
\end{table}

Kiểm định ghép cặp McNemar trên từng đoạn của tập dev cho thấy:
\begin{itemize}
  \item mmBERT-base vượt GuardMerge có ý nghĩa thống kê ($p = \so{dev.mc.guard}$) và nhanh hơn hơn hai lần
        (\so{dev.mmbert.p50}~ms so với \so{dev.guard.p50}~ms). Hướng dùng nền LLM 0,6 tỉ tham số tinh chỉnh LoRA bị
        loại.
  \item Ba bộ mã hóa SafeVi-DualORC-CRF hòa nhau về mặt thống kê trên chỉ số tổng (mmBERT-base so với ViSoBERT:
        $p = \so{dev.mc.visobert}$; so với BamiBERT: $p = \so{dev.mc.bamibert}$).
\end{itemize}

Khi tổng điểm hòa về mặt thống kê, nhóm lựa chọn theo các lát cắt sát nhất với mục tiêu cốt lõi của v2:
\textbf{mmBERT-base} dẫn đầu trên hội thoại nhiều lượt, vượt xa hai bộ mã hóa tiếng Việt trên lát cắt tiếng lóng
(\so{dev.mmbert.tieng} so với \so{dev.visobert.tieng} của ViSoBERT và \so{dev.bamibert.tieng} của BamiBERT), đạt độ đồng
đều vùng miền tốt nhất (chênh lệch chỉ \so{dev.mmbert.gap}), đồng thời nền đa ngôn ngữ giúp tổng quát hóa tốt trước ký
tự Unicode biến dạng (C2). Lợi thế nhẹ của ViSoBERT trên tổng điểm dev chủ yếu nằm ở bình luận một dòng, trùng khớp kho
tiền huấn luyện của nó. Dù trọng số FP32 của mmBERT-base lớn hơn (\so{dev.mmbert.trongso}~MiB), sau lượng tử hóa ONNX
INT8 mô hình chỉ còn \so{tb.int8.mib}~MiB (khoảng 15\% giới hạn 2~GB) và suy luận \so{tb.a.int8.8.p50}~ms trên CPU
laptop; BamiBERT được lưu làm phương án dự phòng siêu nhẹ cho thiết bị cấu hình rất yếu.

\subsection{Đóng góp của từng mô-đun (ablation)}
\label{sec:yc3-ablation}

\paragraph{Đọc toàn bộ đoạn hội thoại.}
Trên \so{ctx.n} đoạn nhiều dòng của tập dev, v2 đạt Macro-F1 \so{ctx.full} khi đọc toàn bộ hội thoại nhưng chỉ còn
\so{ctx.last} khi chỉ đọc dòng cuối (tăng \so{ctx.gain} điểm). Kết quả này khẳng định mô hình thực sự khai thác ngữ cảnh
liên lượt để phát hiện các tình huống bộc lộ dần về cuối đoạn. Một bộ phân loại chỉ dựa vào độ dài đoạn chỉ đạt Macro-F1
\so{ctx.lenonly}, chứng minh v2 không học lối tắt dựa vào số dòng dù bình luận thực tế đều dài một dòng còn hội thoại
tổng hợp phần lớn nhiều dòng.

\paragraph{Bộ chuẩn hóa đầu vào v1.}
Khác với v1 phụ thuộc vào bộ chuẩn hóa 5 tầng, v2 học trực tiếp trên văn bản biến dạng thông qua tăng cường trực tuyến,
phạt nhất quán, FGM và nhánh CNN đọc theo byte. Nhóm đo ảnh hưởng khi bật/tắt bộ chuẩn hóa v1 trên các bộ thử của tập
dev (Bảng~\ref{tab:p4-chuan-hoa}) và xác nhận lại sau khi đóng băng mô hình trên toàn bộ tập kiểm thử
(Bảng~\ref{tab:p4-chuan-hoa-test}).

\begin{table}[H]
\centering
\footnotesize
\setlength{\tabcolsep}{3.5pt}
\caption{Ảnh hưởng khi bật/tắt bộ chuẩn hóa v1 trên các bộ thử của tập dev (Macro-F1, thang 0--100). ``Mức tụt'':
Macro-F1 bản biến dạng trừ bản gốc khi tắt chuẩn hóa; trong ngoặc là chênh lệch khi bật so với khi tắt.}
\label{tab:p4-chuan-hoa}
\begin{tabularx}{\linewidth}{@{}Ycccccc@{}}
\toprule
\textbf{Bộ thử} & \textbf{Đã thấy?} & \textbf{$n$} & \textbf{Bản gốc} & \textbf{Mức tụt} & \textbf{Tắt} & \textbf{Bật} \\
\midrule
Đoạn gốc C0 & -- & \so{cn.cleanc0.n} & -- & -- & \textbf{\so{cn.cleanc0.raw}} & \so{cn.cleanc0.norm} (\so{cn.cleanc0.gain}) \\
Bỏ dấu & Có & \so{cn.accentstrip.n} & \so{cn.accentstrip.goc} & \so{cn.accentstrip.tut} & \textbf{\so{cn.accentstrip.raw}} & \so{cn.accentstrip.norm} (\so{cn.accentstrip.gain}) \\
C1 đã thấy & Có & \so{cn.c1indist.n} & \so{cn.c1indist.goc} & \so{cn.c1indist.tut} & \so{cn.c1indist.raw} & \textbf{\so{cn.c1indist.norm}} (\so{cn.c1indist.gain}) \\
C1 giữ ngoài huấn luyện & Không & \so{cn.c1heldoutcontrolled.n} & \so{cn.c1heldoutcontrolled.goc} & \so{cn.c1heldoutcontrolled.tut} & \textbf{\so{cn.c1heldoutcontrolled.raw}} & \so{cn.c1heldoutcontrolled.norm} (\so{cn.c1heldoutcontrolled.gain}) \\
Chữ đồng dạng (homoglyph) & Không & \so{cn.blindspothomoglyph.n} & \so{cn.blindspothomoglyph.goc} & \so{cn.blindspothomoglyph.tut} & \textbf{\so{cn.blindspothomoglyph.raw}} & \so{cn.blindspothomoglyph.norm} (\so{cn.blindspothomoglyph.gain}) \\
Ký tự vô hình & Không & \so{cn.blindspotzerowidth.n} & \so{cn.blindspotzerowidth.goc} & \so{cn.blindspotzerowidth.tut} & \textbf{\so{cn.blindspotzerowidth.raw}} & \so{cn.blindspotzerowidth.norm} (\so{cn.blindspotzerowidth.gain}) \\
\bottomrule
\end{tabularx}
\end{table}

\begin{table}[H]
\centering
\small
\caption{Kiểm chứng bật/tắt bộ chuẩn hóa v1 trên toàn bộ tập kiểm thử sau khi đóng băng mô hình và ngưỡng (thang 0--100;
trong ngoặc là chênh lệch khi bật so với khi tắt).}
\label{tab:p4-chuan-hoa-test}
\begin{tabular}{lccccc}
\toprule
 & & \multicolumn{2}{c}{\textbf{Macro-F1}} & \multicolumn{2}{c}{\textbf{BC F1 theo ý}} \\
\cmidrule(lr){3-4}\cmidrule(lr){5-6}
\textbf{Tập} & \textbf{$n$} & \textbf{Tắt} & \textbf{Bật} & \textbf{Tắt} & \textbf{Bật} \\
\midrule
C0 & \so{cnt.C0.n} & \textbf{\so{cnt.C0.raw}} & \so{cnt.C0.norm} (\so{cnt.C0.gain}) & \so{cnt.C0.bc.raw} & \textbf{\so{cnt.C0.bc.norm}} (\so{cnt.C0.bc.gain}) \\
C1 & \so{cnt.C1.n} & \so{cnt.C1.raw} & \textbf{\so{cnt.C1.norm}} (\so{cnt.C1.gain}) & \so{cnt.C1.bc.raw} & \textbf{\so{cnt.C1.bc.norm}} (\so{cnt.C1.bc.gain}) \\
C2 & \so{cnt.C2.n} & \textbf{\so{cnt.C2.raw}} & \so{cnt.C2.norm} (\so{cnt.C2.gain}) & \textbf{\so{cnt.C2.bc.raw}} & \so{cnt.C2.bc.norm} (\so{cnt.C2.bc.gain}) \\
Biến dạng chưa thấy khi huấn luyện & \so{cnt.chuathay.n} & \textbf{\so{cnt.chuathay.raw}} & \so{cnt.chuathay.norm} (\so{cnt.chuathay.gain}) & \textbf{\so{cnt.chuathay.bc.raw}} & \so{cnt.chuathay.bc.norm} (\so{cnt.chuathay.bc.gain}) \\
\bottomrule
\end{tabular}
\end{table}

Kết quả trên cả tập dev lẫn tập kiểm thử hoàn toàn nhất quán: bộ chuẩn hóa v1 chỉ hỗ trợ một số biến dạng C1 đã biết
(trên dev \so{cn.c1indist.gain}; trên test C1 chủ yếu nhờ chèn ký tự \so{cnt.kieu.chenkytu.gain}), nhưng lại làm giảm
Macro-F1 trên đoạn gốc C0 (\so{cnt.C0.gain}), nhóm điểm mù C2 (\so{cnt.C2.gain}, như ký tự vô hình
\so{cnt.kieu.zerowidth.gain}, chèn dấu cách \so{cnt.kieu.spaceinjection.gain}) và nhóm biến dạng chưa thấy
(\so{cnt.chuathay.gain}) do đôi khi chuyển đổi sai các từ mô hình vốn đã hiểu đúng. Vì vậy, v2 mặc định \textbf{tắt}
bộ chuẩn hóa v1 và suy luận trực tiếp trên văn bản gốc; năng lực chống né lọc nay nằm trực tiếp trong biểu diễn của mô
hình thay vì bộ luật cứng.

\paragraph{Bộ mã hóa nền.}
Từ Bảng~\ref{tab:p4-ung-vien}, giữ nguyên kiến trúc SafeVi-DualORC-CRF và thay đổi bộ mã hóa nền giữ điểm tổng trong
khoảng hòa thống kê, nhưng tạo khác biệt rõ về khả năng xử lý tiếng lóng, cân bằng vùng miền và chi phí triển khai.

\paragraph{Lượng tử hóa FP32 $\to$ INT8.}
Để tối ưu triển khai cục bộ, mô hình xuất sang ONNX được lượng tử hóa INT8 kiểu lai: bảng embedding nén INT8 đối xứng
theo hàng; các lớp attention và lớp vào MLP lượng tử hóa động cả trọng số lẫn đầu vào; lớp ra MLP chỉ lượng tử hóa trọng
số INT8 (tránh cắt ngọn ngoại lai); đầu phân loại, đầu gắn thẻ và tầng CRF giữ nguyên FP32 để bảo toàn hiệu chuẩn xác
suất. Nhờ đó, dung lượng tệp giảm từ \so{tb.fp32.mib}~MiB xuống \so{tb.int8.mib}~MiB (nén khoảng \so{tb.nen} lần), trong
khi chất lượng trên toàn bộ \so{lt.ALL.n} đoạn kiểm thử gần như bảo toàn tuyệt đối (Bảng~\ref{tab:p4-luong-tu}):
Macro-F1 chỉ chênh \so{lt.ALL.d.macro} điểm (không có ý nghĩa thống kê, McNemar $p = \so{lt.ALL.p}$), BC F1 theo ý thậm
chí tăng nhẹ \so{lt.ALL.d.bcy} điểm, và FPR trên C0, C1 giữ nguyên hoàn toàn.

\begin{table}[H]
\centering
\small
\caption{Đánh đổi chất lượng khi lượng tử hóa ONNX FP32 $\to$ INT8 trên tập kiểm thử (thang 0--100; FPR tính theo \%;
cùng tập và cùng ngưỡng đã đóng băng; trong ngoặc là INT8 trừ FP32).}
\label{tab:p4-luong-tu}
\begin{tabular}{lrcccccc}
\toprule
 & & \multicolumn{2}{c}{\textbf{Macro-F1}} & \multicolumn{2}{c}{\textbf{BC F1 theo ý}} & \multicolumn{2}{c}{\textbf{FPR}} \\
\cmidrule(lr){3-4}\cmidrule(lr){5-6}\cmidrule(lr){7-8}
\textbf{Tập} & \textbf{$n$} & \textbf{FP32} & \textbf{INT8} & \textbf{FP32} & \textbf{INT8} & \textbf{FP32} & \textbf{INT8} \\
\midrule
C0 & \so{lt.C0.n} & \so{lt.C0.fp32.macro} & \so{lt.C0.int8.macro} (\so{lt.C0.d.macro}) & \so{lt.C0.fp32.bcy} & \textbf{\so{lt.C0.int8.bcy}} (\so{lt.C0.d.bcy}) & \so{lt.C0.fp32.fpr} & \so{lt.C0.int8.fpr} \\
C1 & \so{lt.C1.n} & \so{lt.C1.fp32.macro} & \so{lt.C1.int8.macro} (\so{lt.C1.d.macro}) & \so{lt.C1.fp32.bcy} & \textbf{\so{lt.C1.int8.bcy}} (\so{lt.C1.d.bcy}) & \so{lt.C1.fp32.fpr} & \so{lt.C1.int8.fpr} \\
C2 & \so{lt.C2.n} & \so{lt.C2.fp32.macro} & \so{lt.C2.int8.macro} (\so{lt.C2.d.macro}) & \so{lt.C2.fp32.bcy} & \textbf{\so{lt.C2.int8.bcy}} (\so{lt.C2.d.bcy}) & \so{lt.C2.fp32.fpr} & \so{lt.C2.int8.fpr} \\
\midrule
\textbf{Tất cả} & \so{lt.ALL.n} & \so{lt.ALL.fp32.macro} & \so{lt.ALL.int8.macro} (\so{lt.ALL.d.macro}) & \so{lt.ALL.fp32.bcy} & \textbf{\so{lt.ALL.int8.bcy}} (\so{lt.ALL.d.bcy}) & \so{lt.ALL.fp32.fpr} & \so{lt.ALL.int8.fpr} \\
\bottomrule
\end{tabular}
\end{table}

\subsection{Chi phí vận hành trên thiết bị đích}
\label{sec:yc3-chi-phi}

Bảng~\ref{tab:p4-chi-phi} đối chiếu tổng quan chi phí triển khai và vận hành giữa ba hệ thống.

\begin{table}[H]
\centering
\small
\caption{Tổng hợp chi phí vận hành của ba hệ thống. ``--'': chưa đo hoặc không đo được từ phía thiết bị người dùng.}
\label{tab:p4-chi-phi}
\begin{tabularx}{\linewidth}{lYYY}
\toprule
\textbf{Chỉ tiêu} & \textbf{v1 (ViHateT5)} & \textbf{gpt-oss-120b (API)} & \textbf{v2 (mmBERT-base)} \\
\midrule
Môi trường chạy & Trình duyệt (ONNX) & Máy chủ bên ngoài & CPU laptop, điện thoại Android hoặc trình duyệt \\
Dung lượng mô hình & \so{v1.onnx}~MiB (ONNX FP32) & -- & \textbf{\so{tb.int8.mib}~MiB} (ONNX INT8; FP32 \so{v2.trongso}~MiB) \\
Bộ nhớ khi chạy & -- & Không đo được phía client & RAM đỉnh \so{tb.a.int8.8.ram}~MiB (laptop CPU, 8 luồng); \so{tb.mob.4.ram}~MiB (Android, 4 luồng); trình duyệt tăng \so{tb.web.ram.tang}~MiB \\
Độ trễ mỗi đoạn & -- & p50 \so{api.C0.p50}~s, p95 \so{api.C0.p95}~s (\so{api.songsong} yêu cầu song song) & Laptop CPU (8 luồng): \textbf{p50 \so{tb.a.int8.8.p50}~ms, p95 \so{tb.a.int8.8.p95}~ms}; Android (4 luồng): \textbf{p50 \so{tb.mob.4.p50}~ms, p95 \so{tb.mob.4.p95}~ms}; trình duyệt (4 luồng): p50 \so{tb.web.p50}~ms, p95 \so{tb.web.p95}~ms \\
Năng lượng mỗi đoạn & -- & Không đo được (máy chủ) & \so{tb.b.mj}~mJ (CPU, đo pin), \so{tb.web.mj}~mJ (trình duyệt) \\
Chi phí API & 0 & \so{api.C0.usd1000}~USD / 1.000 đoạn & 0 \\
Gửi tin nhắn ra ngoài & Không & Có & Không \\
Chi phí huấn luyện & (vòng loại) & Không huấn luyện & \so{dev.mmbert.trainphut}~phút, đỉnh \so{dev.mmbert.trainvram}~MiB VRAM (\so{v2.gpu}) \\
\bottomrule
\end{tabularx}
\end{table}

\begin{itemize}
  \item \textbf{LLM qua API rẻ nhưng không phù hợp mục tiêu sản phẩm:} dù chi phí mỗi 1.000 đoạn là
        \so{api.C0.usd1000}~USD, phương án API buộc phải gửi nội dung tin nhắn riêng tư ra máy chủ bên ngoài, có độ trễ
        lớn (p50 \so{api.C0.p50}~s, chậm hơn trên một trăm lần so với v2 trên CPU) và phụ thuộc kết nối mạng.
  \item \textbf{Chi phí xây dựng dữ liệu huấn luyện thấp:} toàn bộ quy trình gán nhãn lại, sinh hội thoại và lọc độc lập
        chỉ tốn \so{cp.tong.usd}~USD (Mục~\ref{sec:p1-chiphi}), là chi phí một lần trong giai đoạn phát triển.
\end{itemize}

\paragraph{Đo hiệu năng, năng lượng và độ nhất quán trên thiết bị đích.}
Mô hình ONNX INT8 được đánh giá trực tiếp ở chế độ tương tác \textbf{batch~1} (xử lý từng đoạn khi người dùng nhận tin)
trên ba môi trường đích: CPU laptop (\so{tb.a.cpu}, \texttt{onnxruntime}), điện thoại Android (\so{tb.mob.may},
\texttt{onnxruntime-android}) và trình duyệt (\so{tb.web.trinhduyet}, \so{tb.web.backend}, \texttt{onnxruntime-web}) qua
ứng dụng web v2 (Bảng~\ref{tab:p4-thiet-bi}, Bảng~\ref{tab:p4-nang-luong} và Bảng~\ref{tab:p4-khop}).

\begin{table}[H]
\centering
\small
\setlength{\tabcolsep}{4pt}
\caption{Độ trễ mỗi đoạn (ms, batch~1), RAM đỉnh và thời gian nạp mô hình của v2 trên thiết bị đích.}
\label{tab:p4-thiet-bi}
\begin{tabularx}{\linewidth}{@{}Ycrrrrr@{}}
\toprule
\textbf{Cấu hình} & \textbf{Luồng} & \textbf{$n$} & \textbf{p50 (ms)} & \textbf{p95 (ms)} & \textbf{RAM đỉnh (MiB)} & \textbf{Nạp (s)} \\
\midrule
PyTorch FP32 (CPU laptop) & 8 & \so{tb.a.torch.8.n} & \so{tb.a.torch.8.p50} & \so{tb.a.torch.8.p95} & \so{tb.a.torch.8.ram} & \so{tb.a.torch.8.nap} \\
ONNX FP32 (CPU laptop) & 8 & \so{tb.a.fp32.8.n} & \so{tb.a.fp32.8.p50} & \so{tb.a.fp32.8.p95} & \so{tb.a.fp32.8.ram} & \so{tb.a.fp32.8.nap} \\
\textbf{ONNX INT8 (CPU laptop)} & \textbf{8} & \so{tb.a.int8.8.n} & \textbf{\so{tb.a.int8.8.p50}} & \textbf{\so{tb.a.int8.8.p95}} & \textbf{\so{tb.a.int8.8.ram}} & \textbf{\so{tb.a.int8.8.nap}} \\
ONNX INT8 (CPU laptop) & 4 & \so{tb.a.int8.4.n} & \so{tb.a.int8.4.p50} & \so{tb.a.int8.4.p95} & \so{tb.a.int8.4.ram} & \so{tb.a.int8.4.nap} \\
ONNX INT8 (CPU laptop) & 1 & \so{tb.a.int8.1.n} & \so{tb.a.int8.1.p50} & \so{tb.a.int8.1.p95} & \so{tb.a.int8.1.ram} & \so{tb.a.int8.1.nap} \\
\midrule
\textbf{ONNX INT8 (Android)} & \textbf{4} & \so{tb.mob.4.n} & \textbf{\so{tb.mob.4.p50}} & \textbf{\so{tb.mob.4.p95}} & \textbf{\so{tb.mob.4.ram}}$^{\dagger}$ & \textbf{\so{tb.mob.4.nap}} \\
ONNX INT8 (trình duyệt) & 4 & \so{tb.web.n} & \so{tb.web.p50} & \so{tb.web.p95} & +\so{tb.web.ram.tang}$^{*}$ & \so{tb.web.nap} \\
\bottomrule
\end{tabularx}

\smallskip
\parbox{\linewidth}{\footnotesize $^{\dagger}$ Đo trên \so{tb.mob.may} ở cấu hình tối ưu (Total~PSS thực chiếm
\so{tb.mob.4.pss}~MiB).\\
$^{*}$ Phần RAM trình duyệt tăng thêm khi suy luận so với trang trống (\so{tb.web.ram.trong}~MiB); đỉnh toàn trình duyệt
là \so{tb.web.ram.dinh}~MiB (tạm đạt \so{tb.web.ram.nap}~MiB lúc nạp rồi giảm). Dung lượng tải lần đầu là
\so{tb.web.taive}~MiB, sau đó lưu trong bộ nhớ đệm để chạy ngoại tuyến.}
\end{table}

\begin{table}[H]
\centering
\footnotesize
\setlength{\tabcolsep}{3.5pt}
\caption{Năng lượng tiêu thụ của bản ONNX INT8 đo từ cảm biến xả pin Windows khi chạy liên tục trên laptop (4 luồng).
``Tăng thêm'': công suất đang chạy trừ công suất lúc nghỉ.}
\label{tab:p4-nang-luong}
\begin{tabularx}{\linewidth}{@{}lCCCCC@{}}
\toprule
\textbf{Môi trường} & \textbf{Nghỉ (W)} & \textbf{Đang chạy (W)} & \textbf{mJ/đoạn (tăng thêm)} & \textbf{mJ/đoạn (cả máy)} & \textbf{Wh/1.000 đoạn (tăng thêm)} \\
\midrule
CPU (4 luồng) & \so{tb.b.pnghi} & \so{tb.b.ptai} & \textbf{\so{tb.b.mj}} & \textbf{\so{tb.b.mjtong}} & \textbf{\so{tb.b.wh}} \\
Trình duyệt (4 luồng) & \so{tb.web.pnghi} & \so{tb.web.ptai} & \so{tb.web.mj} & \so{tb.web.mjtong} & \so{tb.web.wh} \\
\bottomrule
\end{tabularx}
\end{table}

\begin{table}[H]
\centering
\small
\caption{Chất lượng suy luận của bản triển khai ONNX INT8 trên toàn bộ \so{tb.int8.n} đoạn kiểm thử (C0, C1 và C2; thang
0--100).}
\label{tab:p4-khop}
\begin{tabularx}{\linewidth}{YCCC}
\toprule
\textbf{Môi trường suy luận} & \textbf{Macro-F1} & \textbf{BC F1 theo ý} & \textbf{Khớp nhãn với FP32} \\
\midrule
PyTorch FP32 (bản đánh giá chính thức) & \so{tb.fp32.macro} & \so{tb.fp32.bcy} & -- \\
ONNX INT8 (CPU, Python) & \so{tb.int8.macro} & \textbf{\so{tb.int8.bcy}} & \so{tb.int8.khop} \\
ONNX INT8 (trình duyệt, WASM 4 luồng) & \so{tb.web.macro} & \so{tb.web.bcy} & \so{tb.web.khop.torch} \\
\bottomrule
\end{tabularx}
\end{table}

Một số kết luận chính về vận hành thực tế:
\begin{itemize}
  \item \textbf{Hiệu năng cao và tiết kiệm tài nguyên trên CPU laptop:} so với PyTorch FP32 (8 luồng), bản ONNX INT8 rút
        ngắn thời gian nạp mô hình hơn 13 lần (\so{tb.a.torch.8.nap}~s $\to$ \so{tb.a.int8.8.nap}~s), giảm RAM đỉnh gần
        5 lần (\so{tb.a.torch.8.ram}~MiB $\to$ \so{tb.a.int8.8.ram}~MiB) và đạt độ trễ p50 chỉ \so{tb.a.int8.8.p50}~ms
        (nhanh hơn cả FP32 trên GPU \so{v2.gpu} ở batch~1 và vượt xa yêu cầu phản hồi thời gian thực).
  \item \textbf{Triển khai trực tiếp trên điện thoại Android đạt hiệu năng tương đương laptop:} trên \so{tb.mob.may} (4
        luồng), bản ONNX INT8 đạt độ trễ p50 \so{tb.mob.4.p50}~ms, p95 \so{tb.mob.4.p95}~ms (nhanh hơn laptop CPU ở cùng
        mức 4 luồng là \so{tb.a.int8.4.p50}~ms), nạp mô hình trong \so{tb.mob.4.nap}~s và RAM đỉnh chỉ
        \so{tb.mob.4.ram}~MiB (Total~PSS \so{tb.mob.4.pss}~MiB, thấp hơn 31\% so với laptop), đồng thời khớp nhãn
        \so{tb.mob.parity.khop} với bản FP32 trên tập đối chiếu \so{tb.mob.parity.n} đoạn C0.
  \item \textbf{Chất lượng ổn định giữa các runtime cục bộ:} trên trọn bộ \so{tb.int8.n} đoạn kiểm thử, bản ONNX INT8
        trên CPU và trình duyệt đạt Macro-F1 và BC F1 gần như tương đương PyTorch FP32 (Bảng~\ref{tab:p4-khop}), khớp
        nhãn lần lượt \so{tb.int8.khop} và \so{tb.web.khop.torch} (chênh lệch nhỏ phát sinh từ sai khác số học của lượng
        tử hóa động ở các đoạn nằm sát ranh giới hai nhãn; riêng trên C0/C1 độ khớp đạt \so{tb.int8.C0.khop} và
        \so{tb.int8.C1.khop}).
  \item \textbf{Suy luận trực tiếp trong trình duyệt và tiêu thụ năng lượng:} bản trình duyệt chạy ngoại tuyến qua
        WebAssembly đạt p50 \so{tb.web.p50}~ms (chậm hơn bản CPU gốc do tập lệnh vector WASM chỉ rộng 128 bit, nhưng vẫn
        đáp ứng tốt khi xử lý tin nhắn đến theo từng lượt). Phép đo năng lượng sử dụng cảm biến pin Windows qua cửa sổ
        chạy liên tục tối thiểu \so{tb.b.giay}~giây kèm hai cửa sổ nghỉ trước--sau, ghi nhận mức tiêu thụ tăng thêm rất
        thấp (\so{tb.b.wh}~Wh trên CPU và \so{tb.web.wh}~Wh trên trình duyệt cho mỗi 1.000 đoạn).
\end{itemize}

\section{Hạn chế và hướng phát triển}
\label{sec:han-che}

\paragraph{Hạn chế hiện tại.}
\begin{itemize}
  \item \textbf{Cỡ mẫu tập kiểm thử còn khiêm tốn} (\so{te.tong} đoạn gốc do thành viên nhóm tự biên soạn), khiến các
        lát cắt nhỏ theo phương ngữ hoặc kênh giao tiếp chịu biên độ dao động lớn hơn.
  \item \textbf{Tỉ lệ báo nhầm trên nhóm kịch bản chưa thấy còn cao,} chủ yếu do ngưỡng phân loại hiện dùng chung cho
        toàn bộ mô hình và thiếu các đoạn đối chứng an toàn tương ứng cho các tình huống ngoài tập huấn luyện.
  \item \textbf{Hiệu năng giảm trên kịch bản hoàn toàn mới} (Mục~\ref{sec:yc2}) và độ trễ trên trình duyệt WebAssembly
        còn cao hơn bản CPU gốc.
\end{itemize}

\paragraph{Hướng phát triển tiếp theo.}
\begin{itemize}
  \item Bổ sung các kịch bản chưa thấy (vay trực tuyến, giả giao hàng, kiểm soát người yêu, chuyển nhầm tiền) cùng tin
        nhắn thực tế đã ẩn danh và các đoạn đối chứng an toàn để nâng cao độ khái quát.
  \item Xây dựng tập kiểm định (validation) độc lập do người gán để hiệu chỉnh ngưỡng quyết định riêng cho từng nhãn
        nguy hại, tối ưu đồng thời FPR và độ phủ phát hiện.
  \item Tăng tốc suy luận trong trình duyệt bằng tăng tốc phần cứng (WebGPU) và đánh giá mở rộng trên các dòng điện thoại
        tầm trung, cấu hình thấp (kết hợp phương án dự phòng BamiBERT).
\end{itemize}
